\subsection{The AI-era literature}\label{sec:lit-modern}

Part II rests on foundations that predate generative AI: the task model, transaction costs, the O-ring, knowledge hierarchies and real options. This subsection sets out the theory and evidence of 2018 to 2026 that bears on each component of the model. It keeps older work only where that work is still the canonical statement, and it says where the modern results support, sharpen or change ours. Working papers and preprints are identified as such.

\paragraph{Tasks and the cost crossover.} The task framework remains the base. Automation displaces labor task by task and new tasks reinstate it \citep{acemoglu2018race,acemoglu2019automation}, task displacement accounts for 50 to 70 percent of the change in the U.S. wage structure since 1980 \citep{acemoglu2022tasks}, and a recent review finds broad support for these predictions \citep{restrepo2024automation}. Our crossover is their automation margin, measured against the cheaper of in-house and outsourced work. Two results change how its value should be read. Automation targets jobs that pay rents above workers' outside options, and dissipating those rents offsets 60 to 90 percent of automation's measured productivity gains \citep{acemoglu2026automation}; the savings we compute are private, and part of them is a transfer. Automation also changes the expertise that the remaining tasks require: removing inexpert tasks raises wages and lowers employment, and removing expert tasks does the reverse \citep{autor2025expertise}. The residual review and sign-off roles of Section~\ref{sec:ext-org} should therefore be priced at an expertise-adjusted wage. Most new work comes from augmenting innovations \citep{autor2024new}. Technical exposure overstates what is worth automating: in computer vision, only 23 percent of the wages paid for exposed tasks would be economical to automate at current costs \citep{svanberg2024beyond}, a working-paper estimate of the gap between our feasible and economical layers (Lemma~\ref{lem:ext-nested}). At the aggregate level, task exposure and cost savings imply modest gains over a decade \citep{acemoglu2025simple}, and scenarios in which every task is eventually automated \citep{korinek2024scenarios} correspond to the limit of Proposition~\ref{prop:ext-agi}.

\paragraph{Make or buy with AI agents.} \citet{coase1937nature} and \citet{williamson1985economic} remain the foundations of the boundary of the firm, and AI now moves that boundary in two directions. AI agents that search, negotiate and transact for their principals lower the cost of using markets, make new market designs feasible and create new frictions such as congestion and price obfuscation \citep{shahidi2025coasean}; related preprints survey economies of interacting agents \citep{hadfield2025economy,rothschild2025agentic}. In our model this is a fall in the coordination share $\kappa_O$, which raises the make-or-buy volume of Proposition~\ref{prop:tasks-window} and favors the vendor. In a Brookings conference draft, \citet{garicano2026vanishing} sets out the opposite force: AI lowers the fixed cost of knowledge, erodes specialists' economies of scale and pulls frequent problems in house. That is a fall in the in-house build cost relative to the vendor's pooled cost, which lowers the make-or-buy volume. The model nests both forces, and which one dominates is an empirical question. Vendor markets are already adjusting. Freelance postings for automation-prone writing and coding fell 21 percent within eight months of ChatGPT's release \citep{demirci2025who}, and exposed freelancers lost employment and earnings \citep{hui2024short}. The outsourced jobs that remained were more complex and better paid, so the vendor's price for the residual, harder part of an activity need not fall.

\paragraph{Augmentation, review and over-reliance.} Field experiments document large gains from augmentation: 14 percent more resolutions for support agents, concentrated among novices \citep{brynjolfsson2025generative}, 40 percent less time on writing tasks \citep{noy2023experimental} and faster coding \citep{peng2023impact}. They also bound those gains. Consultants using AI did worse on tasks outside its ``jagged frontier'' \citep{dellacqua2026navigating}. Across 106 experiments, combinations of humans and AI underperformed the better of the two working alone, especially in decision tasks \citep{vaccaro2024when}. Radiologists given AI predictions underweighted them, so that, in a working paper, the best design assigns each case to the radiologist or to the AI rather than to both \citep{agarwal2023combining}. The opposite bias is just as well documented: radiologists' accuracy fell from about 80 percent to between 20 and 46 percent when a purported AI suggestion was wrong \citep{dratsch2023automation}, and interventions that force deliberation reduce over-reliance at a cost in usability \citep{bucinca2021trust}. These results change two of our assumptions. Review does not catch every error, so Proposition~\ref{prop:tasks-review} should carry a detection rate, below which no review share meets a reliability target. Augmentation carries a quality cost that depends on the task, so it is the natural mode for creation tasks and for tasks in which people outperform the AI, not for decisions in which the AI is more accurate. At scale, gains reach costs slowly. Two years after ChatGPT's release, Danish administrative data rule out effects on earnings or hours above 2 percent, while firms add tasks in AI oversight and integration \citep{humlum2025still}, and adoption is unequal and limited by employer restrictions and training \citep{humlum2025unequal}.

\paragraph{Organization.} \citet{garicano2000hierarchies} is still the canonical model of a knowledge hierarchy. When AI is placed in such a hierarchy, autonomous agents benefit the most knowledgeable people and copilots the least knowledgeable \citep{ide2025artificial}; the AI layer of Proposition~\ref{prop:ext-hierarchy} is the autonomous case, and AI in the solver layer is its natural extension. \citet{agrawal2022power} argue that the larger gains from AI require redesigning systems of decisions, and \citet{agrawal2024artificial} show why: AI predictions raise the variance of decisions, which is costly when decisions interact, so AI is worth more when it is adopted as a designed system. This supports our process complementarity but qualifies its sign condition, since an integration premium can be negative until a coordination mechanism is built. Hierarchies flattened as firms invested in AI, with fewer senior and middle-management roles \citep{babina2025firm}. Generative AI has cut junior and implementation roles in exposed startups \citep{gupta2026generative} and early-career employment in exposed occupations \citep{brynjolfsson2025canaries}, so the residual organization is senior-heavy. Developers with Copilot shift from project management to core work \citep{hoffmann2024generative}, individuals with AI match two-person teams without it \citep{dellacqua2026cybernetic}, and solo founders drive the rise in product launches while teams increasingly dominate the top \citep{kim2026generative}. Both sides of Proposition~\ref{prop:ext-team} thus appear in the data.

\paragraph{Chained errors and agent reliability.} \citet{kremer1993oring} remains the canonical O-ring model. In a working paper, \citet{gans2026oring} show that when tasks are quality complements, deciding automation task by task is incomplete and adoption is discrete and bundled. Propositions~\ref{prop:tasks-complement} and \ref{prop:tasks-aon} reach the same conclusion from handoffs, chained errors and released overhead. Errors compound across the steps of compositional reasoning \citep{dziri2023faith}, agent success on long tasks fits a constant hazard per unit of human time \citep{ord2025is}, and per-step accuracy falls as steps accumulate because models condition on their own earlier errors \citep{sinha2025illusion}. Our capability path uses the logistic fit of \citet{kwa2025measuring}. A constant hazard instead shortens the reliability lag of Proposition~\ref{prop:ext-feasible}(b) at 99 percent from 7.7 to 6.1 doublings, so the engine should carry both forms. Horizons also differ by domain: in a METR research note, horizons on visual computer-use tasks are 40 to 100 times shorter than on software tasks but grow at a similar rate, while self-driving improves much more slowly \citep{kwa2025how}. Consistency matters for high-volume work. A leading agent that succeeds on fewer than half of tasks in a single trial succeeds on all of eight trials for fewer than a quarter of retail tasks \citep{yao2025tau}; the best agent completes 30 percent of realistic office tasks \citep{xu2024theagentcompany}; and multi-agent systems fail through errors of design, coordination and verification \citep{cemri2025why}.

\paragraph{Data and rents.} Data are non-rival \citep{jones2020nonrivalry}. Accumulating them has increasing returns early but diminishing returns in the long run \citep{farboodi2021model,farboodi2023data}, while data-driven markets can tip to a single firm \citep{prufer2021competing,hagiu2023data}. A data advantage is therefore a durable scarce input, in the sense of Proposition~\ref{prop:entry}(c), only if the market tips or the data are continually renewed. Foundation models are concentrated by the economics of scaling \citep{korinek2025concentrating}, and the cost of training frontier models has grown 2.4 times a year since 2016 \citep{cottier2024rising}.

\paragraph{Firm value and intangibles.} Markets capitalized labor substitution at once: firms whose workforces were most exposed to generative AI outperformed by 5 percent in the two weeks after ChatGPT's release, more so if they were rich in data \citep{eisfeldt2023generative}. AI investment raised sales and valuations mainly through product innovation \citep{babina2024artificial}. Within firms, occupations whose tasks are on average more exposed lose employment share, while concentrated exposure offsets part of that loss, and firms that adopt AI (instrumented) grow faster \citep{hampole2025artificial}. Intangible capital scales through non-rivalry \citep{crouzet2022economics,corrado2022intangible}, and market prices now give its capitalization parameters \citep{ewens2025measuring,eisfeldt2022intangible}.

\paragraph{Adoption and timing.} Generative AI spread faster than personal computers or the internet, yet it assisted only 1 to 5 percent of U.S. work hours in late 2024 \citep{bick2026rapid}. Under the original question, use in producing goods or services rose from 3.7 percent of firms in September 2023 to 5.4 percent in February 2024 \citep{bonney2024tracking}. Under the broader question introduced in November 2025 (use in any business function), which is a separate series, 18 percent of firms, or 32 percent weighted by employment, used AI in November 2025 to January 2026 \citep{bonney2026microstructure,allen2026monitoring,census2025btoswording}. Most firms (66 percent) use AI for task augmentation, employment reductions are rare (2 percent), and firm performance is positively associated with the breadth of integration \citep{bonney2026microstructure}. This is the J-curve of \citet{brynjolfsson2021productivity} seen at the level of tasks: most firms sit below the readiness threshold, where only augmentation pays. On timing, \citet{mcdonald1986value} and \citet{dixit1994investment} remain the foundations. Optimal adoption of a technology whose risks are learned over time is gradual \citep{acemoglu2024regulating}, and a preprint obtains our pilot and modularity results in a two-period model \citep{tewari2026pilot}.

\paragraph{AI markets and pricing.} Data on API use show rapid open-source entry, persistent price dispersion, open models about 90 percent cheaper than comparable closed ones \citep{demirer2026emerging,demirer2025emerging}, and a preliminary within-model price elasticity of about $-1.1$ across providers of the same open model \citep{demirer2025emerging}, which measures substitution between providers; the authors report that aggregate elasticities are substantially smaller, so whether spending on inference, and with it the compute share of Part I, falls with its price is open. Menu pricing, market power and multihoming are analyzed by \citet{bergemann2025menu}, \citet{gans2024market} and \citet{fradkin2025demand}. We found no verified study of outcome-based pricing of AI agents.

\paragraph{Ownership.} Employment falls after public-to-private buyouts and rises after private-to-private ones, with large but heterogeneous productivity gains \citep{davis2019heterogenous}. Private equity owners reallocate workers rather than exploit monopsony \citep{herkenhoff2025private}, and after a large hospital buyout margins rose while technology adoption was limited \citep{richards2024managing}. We found no study of AI adoption by ownership regime, so the regime presets for the pace of AI are assumptions. Short-termism goes with lower R\&D growth and, in an estimated model, slows growth \citep{terry2023macro}, a modern counterpart of \citet{stein1989efficient}.

\paragraph{Security and regulation.} The NIST profile for generative AI lists twelve risks, among them information security and the configuration of humans and AI \citep{nist2024artificial}, and NIST's taxonomy of adversarial machine learning covers attacks across the life cycle \citep{vassilev2025adversarial}. Tool-using agents remain vulnerable to prompt injection \citep{greshake2023not,debenedetti2024agentdojo}. Losses from cyberattacks exceed their out-of-pocket costs \citep{kamiya2021risk}, cyber risk is priced in stock returns \citep{florackis2023cybersecurity}, and data on AI incidents are being collected \citep{mcgregor2021preventing}. Expected incident loss should therefore include the loss of firm value, a complement to \citet{gordon2002economics}. The EU's Digital Omnibus postpones the AI Act's high-risk obligations to 2 December 2027 for Annex III systems and 2 August 2028 for Annex I systems \citep{eu2026regulation,eu2024regulation}, and the U.S.\ banking agencies' revised model-risk guidance, which superseded SR 11-7 in April 2026, places generative and agentic AI models outside its scope \citep{fed2026sr262,occ2026bulletin13}.

\paragraph{Quantum computing.} Post-quantum cryptography is a known cost with a date attached. NIST approved FIPS 203, 204 and 205 in August 2024 \citep{nist2024modulekem,nist2024moduledsa,nist2024stateless}, and its final transition guidance deprecates quantum-vulnerable public-key algorithms after 2030 and disallows them after 2035 \citep{moody2026transition}. A 2026 executive order and OMB memorandum require post-quantum key establishment in high-value U.S. federal systems by the end of 2030 and post-quantum signatures by the end of 2031, and they direct a proposed rule that would require federal contractors to comply by the end of 2030 \citep{whitehouse2026securing,omb2026execution}. The United Kingdom sets milestones in 2028, 2031 and 2035 \citep{ncsc2025timelines}, and the EU asks critical infrastructure to migrate by the end of 2030 \citep{nis2025coordinated}. The only quantified cost is about \$7.1 billion for U.S. federal civilian systems over 2025 to 2035, a rough order of magnitude \citep{omb2024report}; we found no credible estimate for firms. Estimates of the resources needed for cryptanalysis have fallen: factoring a 2048-bit RSA key is now estimated to need under one million noisy qubits for under a week \citep{gidney2025how}. For computation itself the evidence is narrower. Quadratic speedups are unlikely to yield an advantage on early fault-tolerant machines \citep{babbush2021focus}, practical advantage requires super-quadratic speedups on problems with little data \citep{hoefler2023disentangling}, exponential speedups for hard optimization problems are unlikely and near-term heuristics lack guarantees \citep{abbas2024challenges}, and an exponential advantage in ground-state chemistry has not been shown \citep{lee2023evaluating}. Error correction has crossed the surface-code threshold \citep{google2025quantum}, and a verifiable advantage is claimed for a physics observable \citep{google2025observation}. Claims of beyond-classical simulation have been met by classical simulations that matched one experiment \citep{kim2023evidence,tindall2024efficient} and challenge another \citep{king2025beyond,mauron2025challenging}. We therefore treat quantum computing as a scenario whose only default effect is the cost of migration.

\paragraph{Part I foundations.} \citet{marris1963model}, \citet{solow1971some} and \citet{odagiri1981theory} remain the canonical treatments of a firm that chooses a point on a frontier between growth and profitability. Modern firm-level evidence supplies the structure of our growth technology: firms grow by acquiring customers through marketing and distribution, at roughly constant markups \citep{fitzgerald2024how,argente2025how,afrouzi2023concentration}. On valuation, cash flows of longer duration earn lower returns \citep{weber2018cash,gormsen2023duration}. If anything, this raises the firm's own growth weight relative to the constant discount rate we assume, and widens its gap with the weights in investors' screens.
